\subsection{$A$-模型}\label{amo}

与基于极化选取的几何量子化不同，膜量子化 \cite{GW} 涉及的是 $M$ 的另一种额外结构的选择，正是它使量子化成为可能。

起点是假定实辛流形 $M$ 有一个具有某些良好性质的\textbf{复化} $Y$。特别地，假定 $Y$ 是复流形，带有反全纯对合 $\tau:Y\to Y$，满足 $\tau^2=1$，且 $M$ 是其不动点集的一个分支。此外，假定 $M$ 的实辛形式 $\omega$ 解析延拓为 $Y$ 上的全纯辛形式 $\Omega$，使 $Y$ 成为复辛流形。我们把“$\Omega$ 是 $\omega$ 到 $Y$ 的解析延拓”理解为 $\Omega$ 限制到 $M$ 上等于 $\omega$：
\be\label{circo}\left.\Omega\right|_M=\omega. \ee
满足此性质的非退化全纯 2-形式在 $M$ 的邻域内自动存在，但我们要求 $\Omega$ 整体地存在于 $Y$ 上。由这些假定可得
\be\label{taupl}\tau^*(\Omega)=\bar\Omega,\ee
特别地
\begin{align}\label{tauaction} \tau^*(\Re\,\Omega)&=\Re\,\Omega \cr \tau^*(\Im\,\Omega)&=-\Im\,\Omega. \end{align}
由于 $M$ 在 $\tau$ 下不变而 $\Im\,\Omega$ 是奇的，还有
\be\label{lagm} \left.\Im\,\Omega\right|_M=0. \ee

一个重要特例是：$M$ 是实代数簇，$Y$ 是相应的复代数簇。我们考虑的所有例子都是这种类型，\cite{GW} 中考虑的例子也是如此。

正如几何量子化，膜量子化涉及一个选择，所得的量子 Hilbert 空间本质上依赖于这个选择。特别地，若 $M$ 有一个具备膜量子化所需性质的复化 $Y$，则 $Y$ 总有许多可能。不能延拓为 $Y$ 的全纯辛同构的 $M$ 的辛同构，会把 $Y$ 映到 $M$ 的另一个可能复化 $Y'$，量子化后导致不等价的 Hilbert 空间。几何量子化中的类似陈述是：不保持给定极化的 $M$ 的辛同构把该极化映到不等价的极化，同样导致不等价的 Hilbert 空间。

对膜量子化，我们要把 $Y$ 视为实辛流形，并研究相应的 $A$-模型。$\Re\,\Omega$ 与 $\Im\,\Omega$ 的任何非平凡实线性组合都可视为 $Y$ 上的实辛形式。然而，为了量子化 $M$，我们要 $M$ 是 Lagrangian 子流形。为此，鉴于方程 \eqref{tauaction}，应取 $Y$ 的辛形式为 $\Im\,\Omega$ 的倍数。由于尚未固定 $\omega$ 或 $\Omega$ 的归一化，就简单地取 $Y$ 的辛形式为 $\Im\,\Omega$。当我们谈论 $Y$ 的 $A$-模型时，总是指取此辛形式的 $A$-模型（可能推广到包含 $A$-模型中的 $B$-场）。

需要对 $Y$ 加一个非平凡条件，才能使命中实辛形式 $\Im\,\Omega$ 的 $Y$ 的 $A$-模型良好定义。例如，给定任何作为某个 $M$ 的候选复化的 $Y$，都可以从 $Y$ 中去掉一对 $\tau$-共轭点而得到另一个候选 $Y'$。这似乎是很不自然的例子，人们会期望把它排除。从物理观点看，应对 $Y$ 施加的一般条件并不很清楚。不过，有一个已知的充分条件，足以覆盖本文考虑的所有例子：$Y$ 容许完备的\textbf{超 Kähler 度量}，它（在马上要解释的意义上）延拓了 $Y$ 的复辛结构。这样的完备超 Kähler 度量保证以 $Y$ 为靶空间的量子 $\sigma$-模型存在，而此 $\sigma$-模型的标准扭转就是所需的 $A$-模型。

这里暂停一下，解释为什么 $Y$ 的 $\sigma$-模型作为二维量子场论真正存在是重要的。在第 \ref{abranes} 节我们将回顾 $Y$ 的 $A$-模型如何与 $Y$ 上全纯函数代数的形变量子化相联系。作为形式幂级数环 $\C[[\hbar]]$（$\hbar$ 为形变参数）上的形变，这个构造是形式的，不需要假定 $Y$ 的量子 $\sigma$-模型真正存在。\footnote{导出 $A$-模型的底层 $\sigma$-模型本身可能是通过量子化某个经典系统构造的，其量子参数通常也叫 $\hbar$。那个底层量子参数未必与 $A$-模型的形变参数有简单关系，后者其实是 $\sigma$-模型的一个模。把这个 $A$-模型形变参数称为 $\hbar$，希望不会引起混淆。无论如何，我们通常不显式引入形变参数，而是用实辛形式 $\omega$ 或复辛形式 $\Omega$ 写公式。要在这样的公式中引入 $\hbar$，可以取一个具体的实或复辛形式 $\omega_0$ 或 $\Omega_0$，写成 $\omega=\omega_0/\hbar$ 或 $\Omega=\Omega_0/\hbar$。我们在公式中显式写出 $\hbar$ 的因子时，往往是在得知可以令 $\hbar=1$ 之前的预备阶段。}但对量子化（而不只是形变量子化）而言，我们要得到一个形变代数，其中 $\hbar$ 可视为复参数，而不只是形式幂级数的自变量。这正是量子 $\sigma$-模型真正存在时得到的东西；此时 $\hbar$ 是 $\sigma$-模型的一个模。也许完整、紫外完备的 $\sigma$-模型的存在超出了本文构造所需，但量子 $\sigma$-模型的存在是我们所知道的充分条件。

超 Kähler 流形有一族复结构，它们与四元数单位 $I,J,K$ 在切空间上的作用相联系，满足通常的四元数关系 $I^2=J^2=K^2=-1$，$IJ=K$。与 $I,J,K$ 相联系的是 Kähler 形式的三元组 $\omega_I$、$\omega_J$、$\omega_K$（其中 $\omega_I$ 关于复结构 $I$ 为 $(1,1)$ 型且为正，$\omega_J,\omega_K$ 类似），以及全纯 2-形式 $\Omega_I=\omega_J+\i\omega_K$ 及 $I,J,K$ 的循环置换。复结构与辛形式满足
\be\label{tono}I^t\omega_J=-\omega_K, ~ I^t\omega_K=\omega_J, \ee
及其循环置换，其中 $I^t$ 是 $I$ 的转置，视为作用在 $Y$ 的 1-形式上的线性变换。$Y$ 的超 Kähler 度量为
\be\label{bono} g=I^t \omega_I=J^t\omega_J= K^t\omega_K.\ee

当我们说 $Y$ 的超 Kähler 结构延拓了它作为 $M$ 的复化的复辛结构时，意思是：$I$ 是 $Y$ 作为 $M$ 的复化时的复结构，相应的全纯辛形式为 $\Omega=\Omega_I=\omega_J+\i \omega_K$。假定 $Y$ 有对合 $\tau$，满足 $\tau^*(g)=g$ 以及先前陈述的条件，这些条件现在成为 $\tau^*(I)=-I$、$\tau^*(\omega_J)=\omega_J$、$\tau^*(\omega_K)=-\omega_K$。由此及方程 \eqref{bono} 可进一步推出
\be\label{tabush}\tau^*\begin{pmatrix} I\cr J\cr K\end{pmatrix}=\begin{pmatrix} -I \cr J \cr -K\end{pmatrix},~~~~
    \tau^*\begin{pmatrix} \omega_I\cr \omega_J\cr \omega_K\end{pmatrix}=\begin{pmatrix} -\omega_I \cr \omega_J \cr -\omega_K\end{pmatrix} .\ee

研究 $Y$ 的 $A$-模型时，我们采用适应此超 Kähler 情形的记号：用 $I$ 表示 $Y$ 作为 $M$ 的复化时的复结构，用 $\omega_J$ 表示 $\Re\,\Omega$，用 $\omega_K$ 表示 $\Im\,\Omega$。

一类重要的例子是：$M$ 为实仿射簇，即它是某个 $\R^N$ 的子流形，由 $\R^N$ 的坐标 $x_1,x_2,\cdots, x_N$ 中某些实值多项式的零点定义。此时可取 $Y$ 为相应的复仿射簇（即 $\C^N$ 中由同样多项式——现在把 $x_i$ 视为复变量——的零点定义的子流形）。

量子化 $M$ 时，可观测量将对应于可解析延拓为 $Y$ 上全纯函数的 $M$ 上函数。但重要的是决定允许哪一类全纯函数：只考虑 $x_i$ 的多项式，还是也允许在无穷远指数增长？由于可观测量将对应 $A$-模型的物理弦态，这个问题有些类似于问允许哪一类量子力学波函数。我们只允许 $Y$ 上的多项式函数，而不允许指数增长函数，因为这样理论简单得多。

\cite{GW} 研究过的几个简单仿射簇例子是 $\SU(2)$ 的余伴随轨道 $M$
\be\label{cosu} x^2+y^2+z^2  = \j^2 \ee
与 $\SL(2,\R)$ 的余伴随轨道 $M'$
\be\label{osuc}-x^2-y^2+z^2=\j^2, \ee
二者都以实变量 $x,y,z$ 及辛形式 $\omega= \d x\d y/z$ 给出。方程 \eqref{cosu} 中 $\j^2>0$；方程 \eqref{osuc} 中 $\j^2$ 的两个符号都有意义。相应的复簇 $Y$（$\SL(2,\C)$ 的余伴随轨道）通过把 $x,y,z$ 取为复值得到。复化之后方程中的符号无关紧要，为确定起见用方程 \eqref{cosu} 定义 $Y$。$M$ 作为 $x,y,z$ 为实的轨迹嵌入 $Y$，$M'$ 则是 $x,y$ 为虚、$z$ 为实的轨迹。$M$ 是由 $\tau:(x,y,z)\to (\bar x,\bar y,\bar z)$ 定义的反全纯对合的不动点集（假定 $\j^2>0$ 使该不动点集存在），$M'$ 是另一个反全纯对合 $\t\tau:(x,y,z)\to (-\bar x, -\bar y, \bar z)$ 的不动点集。把膜量子化方法用于 $Y$，并分别取对合 $\tau$ 与 $\t\tau$，可以重现关于 $\SU(2)$ 与 $\SL(2,\R)$ 表示论的经典结果。其中一部分在 \cite{GW} 中已有描述，我们将在第 \ref{scbq} 节起作进一步评论。
